%============================================================
% constraints.tex — from the constraints to the algorithm (one-page guide)
% Grow the second table as chapters are written.
%============================================================
\chapter{From Constraints to Algorithm}
\label{ch:constraints}

A contest problem tells you its intended complexity through its limits. A judge runs roughly $10^8$ simple operations per second, so before choosing an algorithm, read $n$ (the length), $V$ (the largest value) and $q$ (the number of queries) and look the budget up here.

\section*{Step 1 --- What the limits allow}

\begingroup\small\renewcommand{\arraystretch}{1.3}
\begin{tabularx}{\linewidth}{@{}l l X@{}}
\toprule
\textbf{Limit} & \textbf{Budget} & \textbf{Typical array algorithms} \\
\midrule
$n \le 10$--$12$ & $\bigO{n!}$ & try every permutation / every order \\
$n \le 20$--$25$ & $\bigO{2^n \cdot n}$ & try every subset (subset DP is the DP book) \\
$n \le 40$ & $\bigO{2^{n/2}}$ & meet in the middle (the \textit{Recursion} book) \\
$n \le 500$ & $\bigO{n^3}$ & all subarrays with an $\bigO{n}$ check; fix two rows of a matrix (Ch.~\ref{ch:matrices}) \\
$n \le 5000$ & $\bigO{n^2}$ & all pairs, all subarrays with running state; bitsets push this to $n\approx 3\cdot10^4$ (Ch.~\ref{ch:bitset}) \\
$n \le 10^5$ & $\bigO{n\sqrt n}$ & square root decomposition, Mo's algorithm (Ch.~\ref{ch:sqrt}, \ref{ch:mo}) \\
$n \le 2\cdot10^5$--$10^6$ & $\bigO{n \log n}$ & sorting and binary search (Ch.~\ref{ch:sorting}, \ref{ch:binary_search}); binary search on the answer (the \textit{Binary Search} book) \\
$n \le 10^6$--$10^7$ & $\bigO{n}$ & one scan: prefix sums, hashing, two pointers, windows (Ch.~\ref{ch:prefix_sums}--\ref{ch:sliding_window}); use fast input \\
$n$ or $V$ up to $10^{18}$ & $\bigO{\log}$ & the array is implicit: a formula, binary search on the answer, or matrix powers --- never store it \\
\bottomrule
\end{tabularx}
\endgroup

\begin{insightbox}[Read the values, not only $n$]
Small values change everything: $a_i \le 100$ allows a frequency array per value and loops over the value domain; $a_i \le 10^6$ allows a sieve and harmonic loops over multiples (Ch.~\ref{ch:number_arrays}); OR, AND and GCD over subarrays take only $\bigO{\log V}$ distinct values per right end (Ch.~\ref{ch:bitwise}). Values up to $10^9$ with $n$ up to $2\cdot10^5$ make sums reach $2\cdot10^{14}$: use \texttt{long long}.
\end{insightbox}

\clearpage
\section*{Step 2 --- What the question asks}

\begingroup\small\renewcommand{\arraystretch}{1.3}
\begin{tabularx}{\linewidth}{@{}>{\raggedright\arraybackslash}p{6.2cm} X@{}}
\toprule
\textbf{The problem says \ldots} & \textbf{Try first} \\
\midrule
``maximum / second maximum / count of \ldots'' over the array & one scan with running state (Ch.~\ref{ch:introduction}) \\
``remove / move / compact in place'', ``$\bigO{1}$ extra memory'' & read/write pointers (Ch.~\ref{ch:introduction}) \\
``sum of the subarray $[l, r]$'' for many queries & prefix sums (Ch.~\ref{ch:prefix_sums}) \\
``add $x$ to every element of $[l, r]$'' many times, read at the end & difference array (Ch.~\ref{ch:prefix_sums}) \\
``number of subarrays with sum / remainder / balance $k$'' & prefix sums + hash map (Ch.~\ref{ch:prefix_sums}) \\
``two elements with sum $k$'', ``seen before?'' & hash map of the elements seen so far (Ch.~\ref{ch:hashing}) \\
sorted input, ``pair / triple with sum \ldots'' & two pointers (Ch.~\ref{ch:two_pointers}) \\
``longest / shortest subarray such that \ldots'' with a monotone condition & sliding window (Ch.~\ref{ch:sliding_window}) \\
``minimum total cost to make all equal'', ``arrange in the best order'' & sort, median, exchange argument (Ch.~\ref{ch:sorting}) \\
``first position where \ldots'', counting in a sorted array & binary search (Ch.~\ref{ch:binary_search}); rotated arrays and beyond: the \textit{Binary Search} book \\
``minimise the maximum'', ``maximise the minimum'' & binary search on the answer (the \textit{Binary Search} book) \\
intervals, ``meetings'', ``overlaps'', ``events'' & sort endpoints and sweep (Ch.~\ref{ch:intervals}) \\
values in $[1, n]$, ``missing'', ``duplicate'', $\bigO{1}$ memory & cyclic sort, marking (Ch.~\ref{ch:index_tricks}) \\
``sum over all subarrays / all pairs of \ldots'' & contribution technique (Ch.~\ref{ch:contribution}) \\
XOR, AND, OR in the statement & per-bit decomposition (Ch.~\ref{ch:bitwise}) \\
many offline range queries, no simple inverse & Mo's algorithm (Ch.~\ref{ch:mo}) or a scanline (Ch.~\ref{ch:scanline}) \\
``groups of equal / consecutive values'', ``runs'', ``segments'' & a group pointer over runs (Ch.~\ref{ch:two_pointers}) \\
``cut the array into two / three parts'' & split points from prefix and suffix sums (Ch.~\ref{ch:prefix_sums}) \\
a function that first decreases, then increases & ternary search (the \textit{Binary Search} book) \\
``design a class'', insert / delete / random in $\bigO{1}$, snapshots & array-based data structure design (Ch.~\ref{ch:design}) \\
``Alice and Bob take turns'' on an array & parity, symmetry, first decisive element, Nim sum (Ch.~\ref{ch:games}) \\
the answer is ``always YES'' once $n$ is large & pigeonhole on prefix sums or remainders (Ch.~\ref{ch:number_arrays}) \\
range queries with updates & square root decomposition (Ch.~\ref{ch:sqrt}); Fenwick and segment trees are the \textit{Trees} book \\
window maximum / minimum, next greater element & monotonic deque / stack (the \textit{Stacks, Queues \& Heaps} book) \\
\bottomrule
\end{tabularx}
\endgroup

\begin{warningbox}[Limits are hints, not proofs]
The table gives the \emph{intended} complexity in most problems, not a guarantee. A large $n$ with ``sum of $n$ over all test cases $\le 2\cdot10^5$'' allows per-test linear work but not a per-test array of size $10^6$ cleared each time. Always estimate with the worst test case.
\end{warningbox}
